% !TEX root = ../main.tex
\vspace{-8pt}
Il file \textit{plot\_utils.py} raccoglie le funzioni di visualizzazione per i dati sperimentali. Il modulo definisce anche una piccola palette colorblind-safe riusata dal fit lineare:

\begin{itemize}
    \item \texttt{C\_BAR = "\#4878CF"}
    \item \texttt{C\_MEAN = "\#D65F5F"}
    \item \texttt{C\_BAND\_A = "\#4878CF"}
    \item \texttt{C\_BAND\_B = "\#EE854A"}
\end{itemize}

\section{\texorpdfstring{histogram(x, *, ddof, bins, bin\_width, hist\_range, label, xlabel, ylabel, title, show\_bin\_ticks, tick\_rotation, decimals, show\_mean, show\_band, show\_legend, show\_grid, xlim, ylim, ax, figsize, dpi, save\_path, title\_fontsize, title\_pad, legend\_fontsize, legend\_loc, hist\_alpha, band\_alpha, grid\_alpha, mean\_symbol)}{histogram(x, *, ddof, bins, bin_width, hist_range, label, xlabel, ylabel, title, show_bin_ticks, tick_rotation, decimals, show_mean, show_band, show_legend, show_grid, xlim, ylim, ax, figsize, dpi, save_path, title_fontsize, title_pad, legend_fontsize, legend_loc, hist_alpha, band_alpha, grid_alpha, mean_symbol)}}
\label{sec:istogramma}
\begin{apidesc}
  \item[\textbf{Scopo}]
    Disegnare un istogramma di un campione sperimentale, con media opzionale, banda \(\pm 1\sigma\), legenda, griglia, salvataggio su file e integrazione con assi matplotlib esistenti.
  \item[\textbf{Contratto sugli input}]
    \texttt{x} deve essere un array monodimensionale non vuoto e contenere solo valori finiti. Con il default \texttt{ddof=0} la funzione supporta anche dataset con un solo punto.
  \item[\textbf{Parametri chiave}]\hfill
    \begin{description}[font=\normalfont\ttfamily\bfseries, labelwidth=8em, leftmargin=8.5em, itemsep=1pt, parsep=0pt, topsep=2pt]
      \item[ddof] intero o float (default \texttt{0}) $\to$ gradi di libertà usati per \(\sigma\)
      \item[bins] intero, stringa o bordi dei bin (default \texttt{"auto"})
      \item[bin\_width] float o \texttt{None} $\to$ larghezza fissa dei bin, mutuamente esclusiva con \texttt{bins}
      \item[hist\_range] tupla o \texttt{None} $\to$ intervallo \texttt{(xmin, xmax)} per il range dell'istogramma o la costruzione dei bin
      \item[ax] \texttt{Axes} o \texttt{None} $\to$ asse esistente su cui disegnare
      \item[save\_path] stringa o \texttt{None} $\to$ percorso di salvataggio della figura
    \end{description}
  \item[\textbf{Restituisce}]
    \texttt{(fig, ax)}: tupla con figura e asse del grafico.
  \item[\textbf{Eccezioni}]
    \texttt{ValueError} se \texttt{x} è vuoto o non finito;
    \texttt{ValueError} se \texttt{xlim}, \texttt{ylim} o \texttt{hist\_range} non hanno esattamente 2 elementi;
    \texttt{ValueError} se \texttt{hist\_range} non soddisfa \texttt{xmin < xmax};
    \texttt{ValueError} se \texttt{bin\_width} è non positivo o incompatibile con \texttt{bins};
    \texttt{ValueError} se \texttt{figsize} non è una coppia di valori positivi quando la funzione crea una nuova figura.
\end{apidesc}

\paragraph{Scelta di default su \texttt{ddof}}
La banda \(\pm 1\sigma\) usa la stessa convenzione del resto del package: per default \texttt{ddof=0}, cioè statistica descrittiva di popolazione. Se si desidera la correzione di Bessel, bisogna passare \texttt{ddof=1} esplicitamente.

\paragraph{Uso con assi esistenti}
Se \texttt{ax} è \texttt{None}, la funzione importa localmente \texttt{matplotlib.pyplot} e crea una nuova figura. Se invece viene passato un asse esistente, \texttt{histogram} disegna direttamente su quell'asse e restituisce la figura padre con \texttt{ax.get\_figure()}.

\paragraph{Esempio d'uso}
\begin{minted}{python}
fig, ax = histogram(
    x,
    bins=10,
    title="Distribuzione delle misure",
    xlabel="lunghezza [mm]",
    show_band=True,
)
\end{minted}
